% ============================================================
%  Constans [Hans Brøchner], "The Christian and the Human as Opposites
%  in the Consciousness of Our Time" I--III
%  [Det Kristelige og det Humane som Modsætninger i vor Tids Bevidsthed]
%  Nyt dansk Maanedsskrift, ed. Vilhelm Møller, vol. 2 (April --
%  September 1871): I. pp. 289--299; II. pp. 433--456; III. pp. 492--510.
%  TRANSLATION (English).
%  Translated from transcription.tex (Danish) in this directory, whose
%  header records how its text was established (per-batch word check and
%  a full second reading at the image).  Method: ../../../TRANSLATION-
%  PLAYBOOK.md.  No other edition or translation was consulted.
%
%  ---- REGISTER ----
%  Scholarly, moderately literal, readable English; American spelling.
%  Long periods are broken only where English will not carry them, never
%  across a page marker.
%
%  ---- TERMINOLOGY ----
%  Aligned with this repo's translations of Om det Religiøse (1869) and
%  "Det ubevidste og det bevidste menneskelige Sjæleliv" (1871):
%    det Kristelige / det Humane = the Christian / the Human; det
%      Religiøse = the religious; Religiøsitet = religiousness;
%      Humanitet = humanity; Menneskehed = humanity (Brøchner uses the
%      two near-synonymously: „Menneskehedens Religion“ p. 500,
%      „Menneskehedens Religiøsitet“ p. 508, „Humanitetens Religiøsitet“
%      pp. 505, 509); Menneskekjærlighed =
%      love of mankind; Livsanskuelse = view of life;
%    Tro / Viden = faith / knowledge; Videnskab = science;
%    Virkelighed = actuality, Realitet = reality (as in Om det Religiøse;
%      harmonized across the five batches in the caller);
%    Væsen = essence (Menneskevæsen = the human essence); Aand = spirit;
%      Tilværelse = existence; Forestilling = representation;
%    sædelig / ethisk = moral / ethical; Inderlighed = inwardness;
%      Mægling = mediation; den Enkelte = the single individual;
%      Forsoning = reconciliation; Salighed = blessedness;
%      Selvhengivelse = self-surrender; Gudmennesket = the God-man.
%  Kierkegaard's Øieblikket = The Moment.
%
%  ---- CONVENTIONS ----
%  - Page markers \opage{N} and "% --- p. N ---" at the same joints as in
%    transcription.tex.
%  - The same installment heads, numerals I.-III., section-dividing rules,
%    paragraphs, footnote, \emph spans (letterspacing in the original),
%    signatures and closing rules as the Danish.  Where the Danish spaces
%    only part of a compound, the emphasis falls on the corresponding part
%    of the English (% TR: at the site).
%  - Quotation marks „...“ -> ``...''.  Latin and French as printed, in
%    \textit (p. 294 keeps its roman guillemets: »pagani«).  The
%    signature „Constans.“ is kept.
%
%  ---- DEFECTS AND ERRATUM: DANISH AS PRINTED, ENGLISH CORRECT ----
%  p. 444: the print's „Viden“ is corrected in the volume's own errata
%  (p. 572) to „Veien“; the English translates „Veien“ (the way).  The
%  other printer's errors (listed in the transcription's header) are
%  rendered in their evident sense and logged in % TR: at each site.
%
%  ---- HOW THE TRANSLATION WAS CHECKED (2026-09-27) ----
%  (1) Five translators (pp. 289--299, 433--444, 445--456, 492--501,
%  502--510), each re-reading its English against the Danish sentence by
%  sentence.  (2) Mechanical audit (tr_audit.py): per printed page the
%  same \opage joints, \emph spans, footnote, section numbers, centre
%  blocks and paragraph breaks as transcription.tex, and an English/
%  Danish word ratio within 0.95--1.60.  (3) Terminology harmonized in
%  the caller (Virkelighed = actuality; Religiøsitet = religiousness).
%  (4) Independent review: five readers who translated none of it read
%  every sentence of the Danish against the English and proposed 13
%  corrections, all applied (two reworded).  Sense: p. 434 the relative
%  clause „hvis Princip“ attached to the religion, not the world;
%  p. 455 „til hvilken ... var naaet“ refers to the consciousness, not
%  to humanity; p. 502 „denne ethiske Bestemmelse af det, der ...“ =
%  the determination of the divine, which designates it as spirit;
%  p. 505 the third „paa“ makes giving this existence a content of
%  eternity a third object of the faith.  Terminology: residual
%  reality/real -> actuality/actual for Virkelighed (pp. 296, 455, and a
%  stale % TR: note p. 495); Væsen = essence p. 494; Endeligheds-
%  sondringen = the separation of finitude pp. 509--510.  English:
%  p. 455 'an actuality'; p. 500 uafviselig = inescapable; p. 505 comma
%  so that the will posits the opposition, not the good.  Menneskehed
%  harmonized to 'humanity' (pp. 293, 297).
% ============================================================
\documentclass[12pt,a4paper,openany]{book}
\usepackage[utf8]{inputenc}
\usepackage[T1]{fontenc}
\usepackage{libertinus}
\usepackage{libertinust1math}
\usepackage{textalpha}   % Greek: Διαψάλματα (p. 348), μεσότης (p. 361)
\usepackage{graphicx}
\usepackage[english]{babel}
\usepackage[protrusion=true,expansion=true]{microtype}
\usepackage{setspace}
\usepackage[margin=1.2in]{geometry}
\usepackage{enumitem}
\usepackage{fancyhdr}
\setlength{\headheight}{28pt}
\addtolength{\topmargin}{-13.5pt}

% Footnotes as the printing sets them: *) (one, p. 436).
\renewcommand{\thefootnote}{*)}

% The article's numbered sections: the number alone, centred on its line,
% exactly as in transcription.tex.
\newcommand{\secnum}[1]{\par\begin{center}#1\end{center}\par\nobreak}

\usepackage{hyperref}
\hypersetup{
  pdftitle={The Christian and the Human as Opposites in the Consciousness of Our Time},
  pdfauthor={Hans Brøchner},
  hidelinks
}

\pagestyle{fancy}
\fancyhf{}
\fancyhead[LE,RO]{\thepage}
\fancyhead[RE]{\textit{Constans [H. Brøchner]: The Christian and the Human}}
\fancyhead[LO]{\textit{Nyt dansk Maanedsskrift, vol.\ 2}}

\setstretch{1.3}

\title{\textbf{The Christian and the Human as Opposites\\ in the Consciousness of Our Time}\\[8pt]
  \large [Det Kristelige og det Humane som Modsætninger i vor Tids Bevidsthed] I--III}
\author{Constans [Hans Brøchner]}
\date{\textit{Nyt dansk Maanedsskrift}, ed.\ Vilhelm Møller, vol.~2
  (April--September 1871), pp.~289--299, 433--456, 492--510\\[10pt]
  \small Translated from the Danish}

% ---- original-page markers ----
% \opage{N} marks where printed page N begins, at the same joint as in
% transcription.tex (mid-sentence where the page turns mid-sentence).
\usepackage{marginnote}
\newcommand{\opage}[1]{\marginnote{\footnotesize\textit{[#1]}}}
\newcommand{\apage}[1]{\marginnote{\footnotesize\textit{[#1]}}}

\begin{document}
\maketitle
\thispagestyle{empty}
\clearpage

% --- p. 289 ---
\opage{289}\begin{center}\textbf{\large The Christian and the Human as Opposites in the Consciousness of Our Time.}\end{center}

\secnum{I.}

What comes forth in our time as a struggle between
faith and knowledge is, despite the central significance of this opposition,
nevertheless only a particular expression of a more comprehensive
struggle: the struggle between the Christian and the Human; and in
this struggle Christianity, as the highest and concluding form of the positive
religions, as their purest and most spiritualized expression,
represents in turn something more universal: the
positively religious in its opposition to the free human.
In this last shape the opposition has already for
millennia been pronounced for consciousness: it broke
forth as soon as human thought, bold and strong, had
found its law and its norm in thought itself, as soon as, in
existence --- which it now wished to cognize for itself, not
% TR: the print's „forlaret“ (for-|laret, the k missing at the line break) is a printer's error for „forklaret“ (explained); translated as „forklaret“
to have explained to it through the riddles of myth --- it had found the same
reason that is its own essence.

Already in the antiquity of Greece free thought came into
opposition to the form of the positively religious that was
given in the Greek popular religion. The many gods with
their strife and struggle, with their human passions,
% --- p. 290 ---
\opage{290}with their human destinies, in short, with all their
finitude, became a stumbling block for the thought that sought in the whole of
existence only the One that reason acknowledges, the
all-encompassing, all-governing, infinite, in which it can find
rest. And thought's opposition to the gods of the popular faith
became an opposition that encompassed the \emph{whole} of human life.
As soon as thought had found in itself its own law as the
law of existence, the human will too had to be seen
as governed by this law: in the demand of reason, not in
the command of a foreign will or in the prescriptions of a priesthood,
the norm for human action had to be sought. In
cognition and in will the human spirit then sought the regulating
principle in itself, not by detaching itself from the whole
of existence, but by comprehending itself as a living member of
this whole, and by finding in itself, as indwelling and
akin to it in essence, the divine that was the ground of itself and of
all existence. On this ground it rested
securely, without any secret sense of lacking what thought was combating;
and where Greek philosophy reaches its high point,
in Aristotle, spirit stands in the free clarity of thought, in the
secure consciousness of its unity with the divine reason,
which in its unfolding is the life of the universe, calm and
free over against a faith in which it sees only single glimpses of truth
wrapped in the motley cover of mythic images, in the
involuntary veiling of fantasy or the deliberate veiling of calculating
cleverness.

Under changing forms, but yet in its essence constantly
the same, the opposition between the free human
and the positively religious comes forth in the Christian
era. From the beginning the struggle is that of the new revelation,
of the divine truth, against that wisdom of the world
in which the spiritual is to find a place by God's
wisdom triumphing over man's. Here the Christian is what
contends against the sanctioned, the established, what has been holy
from antiquity, as the wisdom of the world also comes
to include the faith that had originally grown up out of the natural ground of hu%
% --- p. 291 ---
\opage{291}man life, and which is now seen in the same light
as the thought that once combated it. Christianity
fights here, weak in the outward but mighty in the inward, because what
it combats is the outlived, whose time of unfolding is past,
whose root in consciousness has dried up, that which no longer satisfies
the demands of life, because it expresses the human only one-sidedly
and incompletely. It triumphs as it
seems to succumb: from the blood of the martyrs the church rises
up as the church triumphant, whose cross waves above the armies
of the Roman state, whose spiritual power rises up from the imposing
foundation of an earthly power.

In another form the struggle comes forth in the Middle Ages.
Against the entrenched power of the church, represented by a
priesthood spread over the lands, gathered in one apex
in the Vicar of Christ, who in the world's old capital rules
spiritually over ``the city and the world'', the fight is waged not only
from the side of the state, the secular power, but from the side of that
science which the church wishes to make its handmaid
and, as it were, to shut up in a monastery cell, far from nature and
far from the naturally determined life of society. Thought, which willingly
and trustingly subordinated itself to the church, which
strained its strength to solve the task the church had
set it: to join the dogmas together into the unity of an ecclesiastical doctrine,
whose divine truth, to be sure, was not to be
comprehensible to the human spirit, but was yet to be capable of being
maintained by it as something not-impossible; --- this thought
grew imperceptibly under the relation of servitude, and as it grew,
there awoke its consciousness of its own laws, its need
for freedom and its need for actuality. From the object
that, as the unappropriable and incomprehensible, could be shaped only in an outward
manner, thought turned toward the
appropriable content of actuality, which it could draw from a
nature that no longer frightened it back as something foreign,
and which it could fetch up from the human spirit's own
depths. As philosophy, science separated itself from theology,
free thought from bound thought, and the opposition which in the
% --- p. 292 ---
\opage{292}earlier Middle Ages had been able to remain concealed stood, at the close of the Middle Ages,
clear and sharp before consciousness. ``One is
the truth of theology, another that of philosophy'' now became the watchword, and
with that the divorce was accomplished, secular science
emancipated. The church could still brand as heretics, but its
power over minds was no longer the same as when
an emperor, barefoot and in penitent's garb, had humbly to beg a
pope for release from the ban, and when the foremost of thought
had to recant before the councils what thought affirmed
and the church denied.

At the beginning of the modern age the struggle then became a
struggle between emancipated science and the medieval-scholastic
theology, bound by dogma and governed by the uncomprehended
and incomprehensible; and from this weaker
opponent science soon turned against dogma itself.
In the time when antiquity was reborn and the human spirit
turned back with enthusiasm to the classical models in life and
poetry, in art and science, in order to seek in them a content
that it could appropriate and find nourishment in, the struggle was
waged with passion, and in this struggle philosophy got its
martyrs. Soon, however, the opposition assumed a calmer
shape. Science soon no longer needed to fight
for its existence; secured by its own strength, it became a
power of life that compelled recognition, and against which
the divided church had to defend itself. The \emph{whole} of modern
philosophy, in its double form as realistic and idealistic,
stands by its principle in opposition to the positively religious,
as this is represented by Christianity, and
the \emph{whole} of modern natural science is foreign to it. What
Christianity denies --- nature --- is the foundation of modern science,
and human thought, which Christianity
% TR: p. 292 „Synden“ (broken e in the print) read as the word; no translation issue
sees as dimmed and darkened by sin, which has disturbed
its essence, and as led to truth only by revelation,
becomes for modern philosophy a source of truth,
its governing power its own laws, which it finds confirmed
in nature. And by science's opposition to
% --- p. 293 ---
\opage{293}the positively religious the general consciousness of our time
is determined, directly or indirectly. In the whole view of the world and of life
of the modern age the opposition is pronounced
to positive religion's holding fast to a foreign, superhuman
and supernatural law and norm for cognition
and action. For the human consciousness the essentiality of nature
as the foundation of human life stands firm, and it
does not doubt ``the capacity of unaltered, complete human nature
to attain, from out of its own interior, to the truth of cognition
and of action, the source of which the human essence
harbors within itself.'' And this human consciousness, which distances
itself from positive religion, finds in itself the
living source of the ideally religious by going back to the ground of its essence,
to its fundamental relation to the eternal power of existence.

Everywhere in the lands that lead the cultural development
of the modern age we find in our time the opposition between
the Human and the positively religious sharply pronounced and
clearly conscious; and in the many forms, merely heterogeneous
to a superficial glance, in which the Human expresses itself,
it is not difficult to find the pervasive features, and in the
deeper and more serious of these forms there shows itself everywhere a
striving to find, in the very depths of human nature, a
richer and truer substitute for what the positively religious
once provided the spirit. More and more strongly the power of the
Human asserts itself in the course of cultural development, as the wealth
of human nature unfolds ever more fully before
consciousness; more and more it frees itself from the constraining
bonds of the past. Only in the lands that lie outside
the cultural current or at its outer edges does the
old retain to a certain degree its significance for the spirit and its
power over it. Here one can then see the peculiar
phenomenon that precisely this holding fast to the old is interpreted
as a sign that it is these lands in which the
true vital force has been preserved, that it is they who are to
lead the future history of humanity, while
the great cultured peoples are supposed to be outlived and spiritually exhausted.
% --- p. 294 ---
\opage{294}One forgets to draw a historical parallel that lies close at hand.
When the faith of antiquity withered away and a new power of life rose
triumphant, it was not at those points of the old
world where the life of thought was strongest, the current of culture
mightiest, that the old faith put up the longest and toughest resistance.
It was in out-of-the-way regions, in provinces that fed
on old memories, on recollections of vanished greatness;
it was in small towns and villages that the old was preserved,
and the name paganism for the old heathendom
recalls this relation. The same thing repeats itself in our time
with regard to the faith of the Middle Ages; but now people take
pride in belonging to the »\textit{pagani}«.

\begin{center}\rule{1.8cm}{0.4pt}\end{center}

Under all the changing forms that the opposition
between the positively religious and the Human has assumed
from the first Christian era until our days, there is
one phenomenon that must attract attention: it is
the immobility of the Christian view of life beneath its apparent
development. Only to a fleeting glance does
a development show itself in the Christian view of life; what assumes
the appearance of one is only the reflection of the opposition's
alternation, or the amalgamating of the Christian with
elements inconsistently taken up from what is foreign. In
the Human there shows itself a possibility of development as rich as
its essence, infinite as this; in the Christian view of life
there is, and must be, immobility, both because
of its abstraction, of its nature-less spirituality, and because
of its boundness to authority. This boundness
is, essentially regarded, equally great in Protestantism and Catholicism.
Catholicism has an advantage in that tradition
in it is at every moment borne by a living individuality
on whom the currents of the spirit of the age can exert their
influence; but the infallible head of the church himself is nevertheless
so bound by the sum of tradition that piles
% --- p. 295 ---
\opage{295}up behind him that there can be no room for a free development.
It is the same power of tradition, only halted
at a definite point --- at the Scripture, which is itself only
the monument to the tradition of a single Christian age ---
that in Protestantism binds free inquiry and free
development. When the Reformation protested against the tradition of Catholicism
and, in agreement with the general striving
of the Renaissance, of which the ecclesiastical reformation
was only a single, limited expression, took refuge back
in the original, this original was not what is
primitive in the human spirit --- the spirit's freedom in thought and
action --- but what is first in the sense of finite time: Christianity's
first historical monuments; and one took refuge back
in them only in order, for the benefit of the scriptural tradition, to
renounce the right of thought, the autonomy of action, in order
to submit to an outward, abiding norm, more binding than
that of Catholicism.

\begin{center}\rule{1.8cm}{0.4pt}\end{center}

When it is frequently asserted that the general spiritual
movement of our time is going back toward the positively
religious, that there is generally a need making itself felt
in minds to take this up again, this assertion rests
on a confusion. One confuses the religious,
which is a necessary moment in the fully conceived
human, an expression of human nature itself in
its vital connection with the indwelling divine, with
what is expressed in the form of the positive religions. Of the former
that assertion would in essentials be correct; of the latter
it is untrue. It is that religious element --- which the
spiritless materialism of the 18th century and its renewal
in our days have sought to displace --- that there is now a striving, from so
many sides, in all the deeper and more genuine forms of the
human direction of spirit, to reinstate in its right. As
consciousness, as it were, sinks down into its own depths and there
meets with something higher that is the ground of human life and its
% --- p. 296 ---
\opage{296}power of life, it finds again in itself the religious relation; and since
this has been found by way of cognition, the human spirit can
live in it without falling into discord with itself and giving
itself up, and without denying a fundamental activity of its
essence by bringing knowledge as a sacrifice to faith. As
moral consciousness matures, it becomes religious consciousness,
no longer ``religionless morality''. But this religious,
which has its germ in the Human as such, is something
essentially other than the content of positive religion; it is
the universally, ideally religious. Therefore, at the same time as
it is sought, precisely the confessional differences, indeed even the differences
between the various historical religions, become vanishing,
and they must become so, precisely because it is the \emph{universally}
Human in \emph{religious} form that one seeks to draw forth and unfold.
In our days the Protestant and the Catholic, the
Christian and the Jew, indeed the Christian and the Hindu who is touched
by the breath of life of humanity, can meet in a faith that is essentially
one, even if they have not given this common faith a
common expression. And this striving of the age toward the
religious in its liberation from the historical forms --- a
striving that often lacks only full clarity and courage to
acknowledge itself, --- can also be recognized as hidden
beneath the most diverse, apparently wholly heterogeneous phenomena.
It is hidden, as if stunted and unclear, beneath the
attempt of half-speculative, half-Christian theology at a halfway
rationalizing of dogma, at a halfway humanizing
of the Christian, by means of a halfway giving-up of
what is Christian in principle. --- It is hidden in the attempt of ordinary
theology to go behind the historical, confessional
expressions of Christianity in order to find an original,
ideal Christianity that is supposed to have hovered before the Founder's thought
but never to have found a corresponding expression in actuality; for
this ideal, unhistorical Christianity becomes in actuality
only an expression of the abstractly universal element in the positive
religions, in which they coincide with rational faith. ---
It is present, as if hiddenly at work, in the general concep%
% --- p. 297 ---
\opage{297}tion of the positive forms of religion that have come into being in
our days. Even those who feel no rest in the older positive
religions nevertheless regard the new ones with indifference or
mistrust. The relation of the consciousness of our time to positive
religion becomes the same as its relation to miracle.
Like miracle, positive religion is regarded as something
past; even if one does not deny its reality altogether,
one conceives it as something whose coming-to-be belongs only
to the past, has actuality only there, not in the present.
But this conception is only the indirect denial
of its reality altogether. --- And finally that
striving can also be hidden beneath the very adherence to the
absurd positive forms of religion that our time has created;
for frequently it is precisely an unclear seeking after the free
religious that brings people to cast away the earlier positive
forms, if only to find a temporary rest
in a new positive religion.

\begin{center}\rule{1.8cm}{0.4pt}\end{center}

The universally religious, to which the human consciousness
in our time seeks its way back where it conceives itself
more deeply and more completely, this ``religion of humanity''
that is the expression of the human essence's original ideal
religious need, is clarified for consciousness through knowledge,
through true, thoroughgoing cognition, through
philosophical thought. But its relation to this cognition
sets up no abiding distinction between the thinker's
religion and the religion of the multitude. What is grounded in
the universally human must be capable of being formed in universally
accessible forms of consciousness. Not in every form will the
highest truth, which science brings to clarity, be accessible
to all, but in one form or another it will be able to be
so. ``Monopolies of truth for the benefit of a single class of men
would be wrongs against all the others and real
injuries to humanity'' (Lichtenberg). Truth
% --- p. 298 ---
\opage{298}is what unites, what is to become the property of all, the medium of a common
life. And the religious that has its root in
the very ground of the human essence, the free and freed religious,
will, just as it can become the property of all, satisfy
the \emph{whole} human spirit, satisfy the heart's
need too. It will be able to do so because it is the expression of the human essence
itself, and because what is essential for religious consciousness
has the ability to express itself in the most diverse
forms according to the different stages of advancing human
development.

\begin{center}\rule{1.8cm}{0.4pt}\end{center}

For one who with an unbiased and comprehensive glance regards
the religious movement in the consciousness of our time,
a doubt will not easily arise about its real direction
and about its goal for the future. But there lies in the
historical circumstances under which the modern age has developed
much that can make it explicable that the conception of the
phenomena is so frequently confused, where uncertainty with
regard to matters of principle works together with tendencies
that obscure cognition. If the Christian had been
the historically original, if it had come into a \emph{new}
world, and if it alone had been able to fill minds, then
the conception of it could have been sure and easy: its
real nature would, as it were, have forced itself upon cognition.
But when Christianity came into the world,
great periods of culture were already concluded, distinctive fundamental forms
of the Human were already unfolded and exhausted.
And Christianity, as what is foreign to the world, as what
denies nature, has always been able to fill only one domain
of the spirit, a single need, a single form of life, not to encompass
the human essence in its whole compass. It has
constantly had beside it another element, another
principle of life, and since this developed simultaneously with Christianity's
existence and in interaction with the
Christian, the confusion could easily arise that what had only
% --- p. 299 ---
\opage{299}come into being at the same time as the Christian, partly in
struggle against it, could be seen as having proceeded from the very principle of the
Christian, as its rich unfolding. Thereby it became
possible that what is called Christianity could become a phenomenon
composed of such heterogeneous elements that those who
stand within Christianity, but have received the basic content of their
religious consciousness, its points of support, in
this intergrowth with the foreign, could become
uncertain about the boundaries of the Christian against the Human, and
that even the scientifically educated among the representatives of theology
could regard a sharp separation, in principle, of
the specifically Christian, whether it took place from a
standpoint within or outside Christianity, as an
``indefensible misunderstanding'', while they themselves put forward a
conception of the Christian which, in seeking with a fragile
dialectic to annex the Human, made the Christian
into something blurred and unclear, and only concealed its lack of principle
in the fog-clouds of phrases. ---

We shall in what follows attempt to let representatives
of the two conceptions of Christianity contending in our days
have their say: of the one that posits the opposition
between the Christian and the Human, and of the one
that lets the former enclose the latter. And as we then, by way of
history and of the development of concepts, ground our adherence
to the conception that holds fast to the opposition,
we shall set the Christianly religious beside the religiousness
of humanity, which for us stands as the goal of religious consciousness.
For this religiousness we shall then seek an expression, not
exclusively in any single one of its spokesmen, but in
what, in the various nuances in which it has taken
shape, shows itself as something common and pervasive, as
that in which its essential motives become clear.

\begin{flushright}\textbf{Constans.}\end{flushright}
\begin{center}\rule{2.6cm}{0.4pt}\\[-10pt]\rule{2.6cm}{0.4pt}\end{center}

\clearpage
% --- p. 433 ---
\opage{433}\begin{center}\textbf{\large The Christian and the Human as Opposites in the Consciousness of Our Time.}\end{center}

\secnum{II.}

In the conception of the Christian there stand in our day,
within the domain of Christianity itself, two opposites
facing each other, sharply distinguished in principle. The one
has found its classic expression in our literature through an
eminently gifted personality: through \emph{Søren Kierkegaard}.
For him the Christian stands as that which is
unconditionally heterogeneous with everything worldly and with what is
accessible to thought. The center of Christianity, the God-man,
is the absolutely incomprehensible, the unconditioned
paradox, to the acknowledgment of which in faith it is not the
inferences of reason but the despair of the consciousness of sin that brings one.
% TR: p. 433 the Danish reads „Det er det; hvortil“ with a semicolon (logged in the transcription as printed); rendered in the evident sense „Det er det, hvortil“.
It is that to which one attains not through thought but in spite
of thought and through a break with it. Faith
grasps the incomprehensible, and in holding it fast with the
passion of inwardness, it guards, precisely with all its strength
and all its passion, its incomprehensibility against every
attempt of thought to appropriate it. And just as
Christian faith is conceived by Kierkegaard only as having issued from
a break with thought, so too the Christian
life that springs forth from this faith is conceived as having issued from a
% --- p. 434 ---
\opage{434}break with the world. Christianity is the religion of suffering;
its commandment is to die away from the world; only through this
death does it show the way to life. And this negative relation
to the world is not restricted to a single time, as, for example,
to Christianity's first historical appearance. As little
as the unconditioned paradox of faith ever becomes more comprehensible
with time, so little does the life that rests in it ever
become homogeneous with the life of the world. Christianity, which, by
demanding death, the death of selfhood, and the renunciation of everything
worldly outside and within the individual, awakens the resistance of selfishness and
worldliness, must constantly stand equally hostile
over against ``the world''; it must at all times be mocked,
reviled, hated, persecuted. And where history shows a peaceful
relation, that is precisely the sign that Christianity
has lost its peculiar character, that it has been corrupted
by an admixture with the heterogeneous, that it has lowered
the ideal demand, has denied itself. In the Christendom of our time
Kierkegaard therefore does not recognize Christianity; he
sees in it only a deception carried on through many centuries and carefully
veiled; in the church protected by the state
he sees only a caricature of Christianity, whose preachers
live off Christ instead of living for him. He
raises, doubting, the question whether Christianity still
exists on earth, and as, with his ideal of the Christian
before his eyes, he goes back through the history of the church right to its
first age, he carries that doubt with him right up to it
and asks gloomily: have there ever really been Christians?
has the ideal demand of Christianity ever been fully
realized? And this question is consistent; for can
the religion whose principle is \emph{not} to live in the world, to be dead to the world, really find its perfect expression in the world?

In this conception of the Christian there is expressed at once
the sharpest opposition to the Christianity of ``Christendom''
and to the Human. In opposition to the concept of a
Christendom in which the individual is, as it were, a Christian by birth
and continues to be one by virtue of the connection
% --- p. 435 ---
\opage{435}with something objectively subsisting, Kierkegaard accentuates the concept
of ``the single individual''. Face to face he would set the individual,
as a ``single individual'', with God, with the infinite demand of the religious,
in order to let the decision take place in the depth of inwardness, in
the meeting where the divine and the human touch
in the innermost of the self. Every objective way to the God-relation
he cuts off, not only the one that goes through objective
thought, but also the one that goes through the church. The
subjective, personal relation to the truth becomes, according to
man's conditions of existence, the one true relation: in infinite
concern for himself, in the concern that the consciousness of sin
awakens, the individual is, in the subjectivity-relation of faith,
which is the relation of an infinite passionate interest, as a single individual
to win the relation to the divine as it
is revealed as ``the god in time'', as the paradox that
for thought is the absurd, for faith the infinitely certain,
which in every moment, in fear and trembling, is won out of
objective uncertainty.

``The single individual'' is the point of departure, and had to be so
for one who saw in ``Christendom'' a corruption that
comes to light at once in the mixing of the Christian
and the worldly and in the individual's absorption into the ecclesiastical
community, the impersonal universal. But ``the single individual'' is
not, as was quite recently asserted through a conspicuous misunderstanding,
the end point for Kierkegaard.
First the illusion had to be broken, and it had to
be made clear to the individual that its way to God can
go only through the narrow pass of inwardness, not along the highroad
of the objective; only on this condition could the upheaval
in souls on which Kierkegaard staked his life become
actuality. But once the ``single individuals'' had each for himself
entered into the true God-relation, the concept of the congregation,
of the community, regained its validity, because the single individuals
met in the same eternal goal. There came to be a communal relation
for the single individuals, but in such a way that God was the middle term
between the single individual and the community, not the community
% --- p. 436 ---
\opage{436}the middle term between the single individual and God. This relation
is so clearly stated in Kierkegaard's writings that only an urge
to parade a superiority by getting Kierkegaard's
standpoint subordinated as a moment in a higher ethical-Christian-ecclesiastical-human
one can explain the misunderstanding that ascribes to
Kierkegaard another conception.

But the life that is set as a task for the single individual
while he stands alone face to face with God, and while
he finds in his God-relation the relation to the other individuals,
becomes a life that sharply separates itself from the Human.
What sustains it is that which is inaccessible to thought:
the unconditioned paradox, that which for thought is the absurd.
But the paradox, which conflicts with all the laws and
forms of thought in that it asserts the union of God and the individual
human being in the person of Christ, the coming into being of the eternal in a moment
of time, is for the believer the highest truth. The
true then judges that which conflicts with it as the
untrue, and there remains in the spirit no place for a life of
thought; there remains no validity, not even a relative one,
for thought. The paradox has rightly been designated
``the miracle in the world of thought'', and it has rightly been
emphasized that, just as the miracle in the world of nature cancels
the concept of nature and of natural law, so the miracle in
the world of thought must overturn the laws of thought and let all the
certainty of knowledge dissolve into doubt. The demand on the believing
individual would consistently have to be ``that in its life of thought it is to
cast away all the world's content of actuality merely in order to be with
the paradox of the consciousness of sin in the paradoxical God-relation.''
This must become the demand, because the paradox contradicts not merely
a single form of human thought but \emph{all} human
thought; because here there is not a conflict between two
forms of the rational, but between the rational in its
entirety and the antirational. The absurd, Kierkegaard's
categorical determination for the paradox, excludes \emph{all} thought.\footnote{Cf. Nielsen: Religionsfilosofi, pp. 8, 34, 301 ff., and elsewhere.}

% --- p. 437 ---
\opage{437}But just as the believer's life, as Kierkegaard
conceives it, forms the opposite of the human life
by denying thought, which expresses an essential determination
in the human, so it comes into conflict with that life
because it consistently cannot make room for the concretely
moral, for all those forms of an ethical life that rest on
a natural ground in the human.

According to Kierkegaard's conception the religious relation becomes
not merely the highest relation, to which the ethical relation
points, in pointing beyond itself, but it becomes the
absolute relation in the sense that it cannot ``exhaust
itself or become concrete'' in the ethical relative relations. It
becomes the absolute relation, then, in the sense that it lays claim
for itself to all of man's strength and lays claim to it
in every moment; for only as I always, in every moment, ``relate
myself absolutely to the absolute'' do I preserve the relation
according to its essential determinacy. But then no strength remains
over for the ethical relations: they must vanish before the absolute,
the infinite, precisely because they fall outside it; just as
the finite strength of the creature must vanish when it is placed
over against the infinite might of omnipotence. Therefore the demand is
also made of man that, for the sake of the
infinite religious relation, he be able to renounce \emph{all} the ethical
relations. The religious does not become the soul in the ethical,
does not permeate the ethical relations in such a way that, in
raising them up into a higher sphere of infinity, it itself gains
living actuality through them; rather it stands over against them as
the infinite that consumes the finite, as the divine
that brings death to the human whose eye has seen it.
In Christianity's demand to die away, to hate one's natural
life, it is implied that the content of natural existence
vanishes as a nothing in its unconditioned opposition
to the eternal truth, and with this natural foundation there vanishes
the condition for that unfolding of the ethical through
which alone it can gain actuality. The more decidedly Kierkegaard
stamps his standpoint in the course of his opposition to the
% --- p. 438 ---
\opage{438}unprincipledness of the established order, the more decidedly this
consequence too comes forward. In the polemic against marriage in ``The Moment''
the consequence is stated undisguisedly.

\begin{center}\rule{1.8cm}{0.4pt}\end{center}

The contrast to the Kierkegaardian conception of the Christian
is formed, now, by the one that extends the name of the Christian to
all those forms of spiritual life that have developed simultaneously
with and under the influence of Christianity. It
sees Christianity not as the opposite of the spiritual life of our time
but as its animating principle; it speaks
of a Christian state, a Christian ethics, a Christian art,
a Christian science; in a word: it applies the name of Christianity
to all forms of the Human; it calls
the world in which Christianity is like a foreign
guest a Christian world. The sharp distinction between
the principle of Christianity and that of the world it allows, although it
constantly points to the apostolic as the orienting,
to have been justified only in Christianity's first age, in which
the struggle with the forms of the old world had to lead to an
exclusive relation. From the moment when Christianity became
victorious, on the other hand, its task was not to be to separate
itself from the world, to die away from it, but to permeate
it; its divine power is to show itself precisely in this, that
it has been able to do so, and its capacity, as principle of life,
to permeate the heterogeneous is sought in this, that Christianity
is a new creation, which acts back in an elevating, ``potentiating'' way
on nature and on the spiritual forces and
laws. Not flight from the world but immersion in it, not
struggle against it but peace with it and dominion over
it then becomes Christianity's goal, which in the Christian state
is reached in the way in which it can be reached under the conditions of finitude,
and the state's recognition of the Christian, its favoring
of it in the way of worldly goods, becomes a new
proof that the human Christianity, reconciled with the world,
is the true one. And the same recognition is also rendered
% --- p. 439 ---
\opage{439}to Christianity by reason. Even if the unregenerate, darkened
reason does not comprehend the mysteries of Christianity,
the believing reason, potentiated by the new principle of life, is nonetheless
able to penetrate into their depths, and even if it does not
attain logical or mathematical clarity, it can nonetheless explore
the connection of the articles of faith, and thereby their
ground, and develop a science of faith that is the
summit and norm of all science. Christian philosophy rises, as
thought borrows wings from imagination, to descry in the God of Christianity
an eternal nature, cleansed of all the imperfection of the grosser nature
defiled by man's sin.
It sees in this eternal nature in God the possibility of
asserting the relative justification of the natural in the human,
the possibility of a reconciliation of the oppositions in human life. And
in the union of the divine with the human in
Christ it sees the highest expression of the reconciliation that thought
demands: between the infinite spirit and the finite, between
the creator and the created, spiritual-natural existence. Thus
Christianity is seen as reconciled with the world and with
thought, as their new, consummating principle of life.

For this conception the Kierkegaardian one becomes a damnable
one-sidedness, just as much as the ``emancipated''
reason's assertion of the Human as the true in its
distinction from the Christian. In place of distinction there steps
everywhere mediation; the rich ``higher unity'' asserts itself with
dignified superiority over against the poor one-sidednesses;
the Human is condescendingly taken up as a serving moment,
and with rhetorical pathos the triumphs of Christianity are proclaimed.

To a fleeting glance all this can be very beautiful to
behold, and the need for a reconciliation won without the struggle of the will
and the exertion of thought, a need so widely spread, makes
it psychologically easy to explain that such a theory of leveling
finds approval with the great multitude, which needs not
clarity but reassurance, and which more readily receives the
smooth phrase than the sharp thought. If one looks a little more closely
at the matter, however, one discovers without difficulty that the
% --- p. 440 ---
\opage{440}smooth phrase only binds together loosely stitched oppositions,
that the higher, richer unity that is supposed to have been won
has been procured only by a very poor method, stereotypically
repeated everywhere, which can yield as its result only semblance, not
actuality, phrases, not concepts. The stereotyped method
is simply this: one takes the opposites one wants
to bring together and joins them in one sentence as subject
and predicate, and then one sends the phrase so formed
out into the world with the endorsement that it ``must be thought'', which
presumably is to be understood rather like the ``may be printed''
% TR: p. 440 Danish „maa“ means both "must" and "may"; the article plays on „maa tænkes“ (must / may be thought) against the censor's „maa trykkes“ (may be printed), and the next sentence takes „maa tænkes“ as a mere permission. English cannot keep the one word; rendered "must be thought" throughout and "may be printed" here.
(\textit{imprimatur}) that the censor in his day put on the books that
were sent out into the world. At any rate one cannot well attain to more than a permission to think
the phrase; an actual thinking of
it becomes a sheer impossibility. Since both opposites
find their expression in the phrase, these expressions of theirs must, as it were,
neutralize each other: the predicate that is joined to the heterogeneous
subject must, instead of determining the subject, rob it of
its content and leave behind for consciousness only an indeterminate, vague representation,
a representation of a nothing; and when
the opposites in their determinacy then press forward again
to consciousness, the contradiction and the discord are there anew, and
the unity must vanish, unless consciousness can be lulled into
a gentle slumber by the monotonous music of the phrase. A single
example will suffice to illustrate what has been said. The opposition
between positive religion and philosophy in the conception
of the relation of the divine to the being of existence finds its
pregnant expression in the opposition between the representation of creation
and the concept of a world-unfolding, a cosmogony. For
religion existence is a work of the incomprehensible arbitrariness
of a creating omnipotence; for philosophy existence has being
in virtue of the very eternal necessity of the essence of the divine; it
is an unfolding of the essential determinations of the divine itself.
Here sharp opposites, irreconcilable as it would seem, stand
over against each other; but that method unites
them easily --- in a phrase that ``must be thought''. For first
the representation of creation is rationalized by the concept of cosmogony:
% --- p. 441 ---
\opage{441}it is not to be a God determined as nature-less spirit, who
by an incomprehensible act of omnipotence produces out of nothing
an existence that, as material, has no point of attachment
in the essence of the divine and therefore cannot be comprehended
as its effect. It is to be a God who has a
nature within him, who creates out of free necessity, and it is to
lie in the concept of creation that God produces not something dead
but something living, that he produces a creature that at the same time
produces and develops itself in a given independence.
Therefore creation \emph{must} be \emph{thought} to ground cosmogony,
or the world's self-development, its genesis. But when
this step toward naturalism has been taken, when that which
for the religious consciousness makes creation creation has been
effaced, then the foot is cautiously drawn back, and the text
continues: ``The world's cosmogonic or natural beginning
is the relative, the finite one, which is therefore also split into
sporadic multiplicity. All organization forms itself sporadically,
and viewed from the cosmogonic or the natural
standpoint the world can be said to have countless beginnings,
to begin from infinitely many seeds of life, which, insofar
as the natural standpoint is exclusively held fast,
have their common unity only in chaos. But the infinitely
many natural beginnings are all grounded in the one
creative, supernatural beginning, in the divine Logos-will,
which encloses within itself the wellspring of life and of light and out of its
spiritual depth releases the whole multiplicity of living forces
to self-movement.'' So: creation was cosmogonic,
the supernatural natural, and now in return
cosmogony becomes creation, the natural supernatural. Thus
the unity must be thought; indeed, it would be an indefensible
one-sidedness to separate the determinations. But just as the
creation that was cosmogonic lost, precisely through this predicate,
that which was to make it creation, so now the
cosmogony which, as \emph{grounded} in creation, itself becomes
creation --- in that the relative, natural beginning vanishes
in the supernatural, the material in the ``spiritual depth'' --- loses in
% --- p. 442 ---
\opage{442}return precisely that which makes it cosmogony. We are
then left with two empty representations, two names without any content,
which we are free to think, if indeed we can
think anything by what is contentless; and if we cannot,
and, so as not to let our thought lose itself in the empty, we call
back the content of the concepts and then try to use the permission
to think their unity, the opposites split
apart again. Then creation again stands against cosmogony, arbitrariness
against necessity, imagination's picture of nature
in God's ``spiritual depth'' against real, actual, material
nature, the miracle against the lawful
development of nature valid in its essence, the supernatural against the natural. And the contradiction
is not removed by such logogriphs as that God
``restricts'' his omnipotence but ``does not give it away'', that he
``produces a creature'' that at the same time ``produces itself'',
a creature that possesses a ``given independence''. The unity
evaporates in the empty turns of phrase: both opposites
are distorted, and sense can be brought into the phrases only
if, in place of the half rational, the half
natural, one puts the wholly antirational, the wholly supranaturalist,
takes what belongs to the one opposite
as a mere turn of phrase that is not meant to signify anything further, and lays
all the emphasis on what belongs to the other. But then
mediation has been given up, and we have returned to the conception
that was condemned as a one-sidedness.

\begin{center}\rule{1.8cm}{0.4pt}\end{center}

We could restrict ourselves to demonstrating the frailty of
the method of the theory of mediation; but in order to provide complete
material for the judgment on the \textit{Juste-milieu} theory, we
shall also take the path of history and of the development of concepts, in order along it
to secure for ourselves the understanding of the Christian, the conception of
its essential and principled relation to the Human.

When we would find the historical presuppositions for
the coming into being of Christianity, they are to be sought essentially in the development
of the Greek and of the Jewish spirit. From different
% --- p. 443 ---
\opage{443}points of departure these national spirits, the representatives
of the Aryan and of the Semitic stock, strive toward one goal of development,
and as they meet, the conditions are given for
the coming into being of Christianity. Whatever may have been taken up into it
from other, oriental, elements has in any case
come to it through the Greek or the Jewish.

If we then consider first the development of the \emph{Greek} spirit,
what becomes of essential significance here is the breaking forth of dualism
in Greek thought and view of life. The Greek people,
the representative of the life of beauty in world history, had found
its highest expression for the intuition of existence in the conception
of it as a cosmos, as a world-artwork; it had
found its highest expression for the ideal human life in the conception
of it as a moral artwork, as a harmonious
forming of the natural side in man by reason. The intuition of beauty
was thus for the Greeks the highest; it
poured itself like a radiance over existence and human life,
which was seen as resting in the repose of harmony. And in the relation
to the gods the same intuition of beauty stepped
into the foreground. When a Greek stood before Fidias's
image of the Olympian Zeus or of Pallas Athene, then
his devotion was this: to let the harmony of beauty of the ideal image
% TR: p. 443 Danish as printed „Stihed“ (printer's error for „Stilhed“); rendered "stillness".
sink down over the soul, to bring stillness, repose and
harmony into its desire and passion, to let all pain
be soothed by the intuition of the divine, present
beauty.

In this intuition of existence as beauty the dominion of the ideal
is asserted; nature is subjected to the power of spirit.
But the unity of spirit and nature in the harmony of beauty is nonetheless
only a deceptive unity, deceptive because the artwork, in its
immediate unity for intuition, joins two opposites:
material and form, which stand as originally distinct,
as not sprung from each other. When Greek thought
reaches its full development, when it, as it were, unwraps the idea
from the concealing veil of the sense-world and would comprehend its essence
and its relation to the material that through it becomes a
% --- p. 444 ---
\opage{444}world-artwork, then the opposites separate for thought,
and Greek thought is not able to lead them back from the separation
to a unity firmer and surer than the beauty-unity
of harmony, because, as Greek thought, it has its limit
in the intuition of beauty, is bound by it as by an
insurmountable barrier. When Plato and Aristotle had brought
Greek thought to its highest development, the dualism became
visible, and the period of dissolution of the intuition of beauty begins.
The consciousness of the discord then leads first to attempts
to see the whole of existence in the light of one of the opposites,
and as these fail, it forces forth dualistic theories in
which a \emph{purely spiritual} God receives as his original
opposite dark matter. But the dualistic distinction
between spirit and nature, which tore apart the unity of beauty,
had to contribute to calling forth skepticism; it had to
give spirit, which was separated by a gaping chasm from the
material, mistrust of its cognizing, of its theoretical faculty.
And as this doubt was awakened, it brought decisive consequences with it.
Whereas before, spirit, trusting securely in the power
of its thought, had set that thought the highest goal and had found rest
in its cognition of the divine and of existence,
it now, in mistrust of thought, seeks other sources of
truth, and it believes it finds them in the revelations of the hidden primal being
through the positive religions or through
demonic intermediate beings that connect the inaccessible
deity with finite existence. And with mistrust
of knowledge there had also to follow mistrust of man's moral
strength. For the Greek, knowledge was man's essential determination;
therefore that double mistrust was inseparable.
Recourse was then had once more to the divine and to the intermediaries
that connected it with man, in order to find aid toward
what spirit could not reach by its own strength. Here too
% TR: p. 444 the print reads „Viden“ (knowledge), a misprint corrected in the volume's own errata (p. 572): „Viden“ for „Veien“; translated from „Veien“ (the way).
the way turned aside from philosophy to religion, and it led
through purifications and cleansings, through flight from the
sensuous to the divine being, from which every determination
of matter was excluded. The truth of existence
% --- p. 445 ---
\opage{445}thus lay in the purely spiritual, which rose to an infinite
distance above the imperfection of matter, in which the germ of
evil was sought; and the striving of the human spirit, which in
the classical age of Greece had been satisfied within the
whole of beauty that mirrored the Idea, now went out beyond the world
to a world beyond, where all was spirit.

With this development in the life of the spirit there then followed a consequence
that became of essential significance for posterity. As
everything sensuous, everything natural, lost its significance, the natural barriers
that separated the peoples also lost their significance;
the particularist spirit that divided Greeks from barbarians, and
that saw the true representative of the human only in the free-born Greek,
was replaced by a universalism that recognized
man \emph{as} man, and asserted the human
through the equality of men.
% TR: p. 445 „Menneske-Ligheden“ rendered "the equality of men" (Lighed = equality, also likeness).

What was reached along this path gained strength from the external
political conditions. The Roman world empire broke through
the barriers of the nationalities, crushed the distinctive character of the peoples.
Through suffering, spirits were led away from unsatisfying
finitude to the beyond to which the religions
pointed. But the popular religions and their individually distinct
divine figures were already in the process of melting
together into an unclear unity, and this process of fusion
was favored by the political unity. But while
a Pantheon gathered the many gods and theology ordered
their hosts, there rose above them a highest unity,
lifted above all being and all thought and yet seen as the ground of all
being and all thought; and for this invisible, purely
spiritual divine being the many gods became serving demons,
which led the yearning human spirit to its source, into which
it sank down.

At the point to which we have followed the Greek
spirit in its development, where it had exchanged the harmony of beauty
for spiritualism, national particularism
for universalism, it meets the \emph{Jewish} spirit.
This spirit had, in the time of the old Hebraism, rested securely in the
% --- p. 446 ---
\opage{446}present and this-worldly. The chosen people's contractual relation
to Jehovah, the law with its rewards and punishments,
the sacrificial service and the cultus pointed to a childlike stage of development,
where the spirit has its norm outside itself and its goal
in something finite. Through the prophets the religious was inwardized;
the tribulations of the time, which they interpreted from a religious
standpoint, led out beyond the present to a hope for the future
of a messianic kingdom; the first intimations of an immortality
glimmered forth, and the narrow barriers of national consciousness
began to open. After the exile we then find, in a certain
sense, a regression, an unfree clinging to the traditional.
But on the other hand we find preparations
for a progressive development. We find such in the
representation, formed under Persian influence, of a
spirit-realm that is the archetype and truth of the earthly, in the representation
of a supernatural Messiah, and in the doctrine of the resurrection.
Finally, in the time immediately before the appearance of Christianity,
we see what had been prepared come forth. Through contact with the
Greek, the dualistic distinction now develops more definitely
between a kingdom of the world and of matter, the kingdom of the evil powers,
and a longingly awaited heavenly kingdom in a natureless
actuality. The heavenly kingdom is striven toward
through an ascetic flight from sensuousness, and faith in immortality
takes the form of a hope for a complete
liberation from earthly matter. Connected with this
is a moral inwardness and a universal love of mankind, which
in practice led to the condemnation of slavery (Essenes and Therapeutae).
The representation of the Messiah is developed toward the supernatural.

In the conception of the divine in an abstractly spiritualistic
manner; in the dualistic opposition between the
spiritually divine and the material; in the human spirit's
deep need to find peace and rest in a kingdom that is not
of this world, by a dying away from everything natural,
by flight from the world; in the distrust of being able, by its own power,
to reach this goal --- thus, at the beginning of the Christian
era, the two national spirits meet: the spirit whose fundamental
% --- p. 447 ---
\opage{447}thought was \emph{God}, and the one whose fundamental thought was \emph{man};
and Christianity, which unites the two thoughts in the thought of the \emph{God-man},
is determined in its basic character by its
two essential factors.

Before we demonstrate this in the principal doctrines of Christianity,
we will briefly see the \emph{conceptual
development of the positive religions} lead up to the point where the converging
lines of development of the two national spirits have proved to
meet.

As the human spirit in positive religion, driven by
the need of its essence, seeks a relation to the infinite,
this striving of it toward the infinite is realized only under conditions of finitude:
through a progressive series of forms
the religious consciousness must approach its ideal. The
infinite that it seeks it finds first in nature: nature religion
is the first form of positive religion, nature
the first divine being. But the seeking spirit, which
is not yet able to distinguish itself as something higher from
nature, is not able either to distinguish nature from itself, in the way
that it immediately senses itself as living,
feeling and willing. Nature, before which man bows
in worship, is not conceived \emph{as} nature, but is
humanized, is conceived as a personal, living,
humanlike being. Thus the child sees external things
as living and ensouled; thus it caresses them or
is angry with them, because it ascribes will and thought to them.
As the natural being that is worshipped is humanized,
it becomes, through this very personification, in a certain sense a
% TR: p. 447 only „over“ of „overnaturligt“ is spaced; emphasis on "super" of "supernatural".
\emph{super}natural being. The natural object becomes as it were doubled:
above the sensuous natural thing rises the thing's
genius, its Fylgie, its Ferver, --- to recall the
expressions of the different religions. As the dryad separates from
the tree and yet is the tree, as the naiad separates from the stream
and yet is the stream, so the doubling takes place everywhere, in
the particular and in the comprehensive. What in the doubling stands
as it were behind the natural is then determined as something invis%
% --- p. 448 ---
\opage{448}ible, as something inaccessible to the senses --- thus
nature religion avoids the contradiction that is in it between
representation and actuality, imagination and truth. The contradiction
between the two objects into which the one natural object
has divided is thereby to be canceled, in that the object that is religion's
proper object is withdrawn from the immediate
relation and becomes a being that is object only for faith,
for the spirit. Herewith the conception of the divine in the religion of spirit
is already intimated. And the religion of spirit replaces
nature religion where human development has advanced
so far that man, from a merely naturally determined being,
has developed into a being that distinguishes itself from nature and
as it were concentrates itself within itself. When man, through
human social life, becomes a being that is determined
by conscious laws of reason, when powers other than
natural powers --- political, moral powers --- become the essential
objects of his feeling of dependence, then the world and
nature no longer appear to him as God, but as a being determined by
a divine understanding and a divine will
and dependent on it. As man himself, by will
and understanding, rises above nature and thus becomes a
supranaturalistic being, his God at the same time also becomes
a supranaturalistic being, no longer one with nature,
but lord over nature; and he must finally, in order
to be lord in the strict sense, become creator of nature.
And the many gods that the religions of spirit receive from
the nature religions must consistently melt together into a divine
unity. The unity of human consciousness and spirit,
which comprehends the whole multiplicity of the external world,
is what leads to positing the unity of God.

Thus positive religion works its way up from
the lowest stage of the nature religions to the highest of the religions of spirit,
to the representation of the one, spiritual, creating God. But
precisely when it reaches this point, the divine is
conceived abstractly spiritualistically as the \emph{natureless spiritual}.
The organ of positive religion is not compre%
% --- p. 449 ---
\opage{449}hending thought, but representation and feeling. What it possesses
concentrated in the rich inwardness of feeling it
unfolds in representation; only in this unfolding does
the content of feeling attain definiteness and clarity. But the distinctive mark
% TR: p. 449 the print's „sandelige“ rendered as „sandselige“ (sensuous).
of representation is that, from the sensuous intuition which
it is to carry over into thought, it has retained a form of externality
which brings it about that its content steps out into a finite
distinction and is fixed for the inner eye in this distinction.
Therefore the divine, the source and
ground of all existence, is placed as divine personality above and outside
% TR: p. 449 again only „over“ of „overnaturligt“ spaced; emphasis on "super".
this existence, determined as a \emph{super}natural being, in the
sense that this supernatural forms the \emph{opposition}
to the positive determinations of nature. The \emph{representational distinction}
of the divine from nature as the supernatural
includes the conception of it as the natureless spiritual:
the transcendence of representation includes abstract spiritualism.
But if the conception of the divine as
the one, spiritual, creating God is the goal of the positive religions,
then abstract spiritualism must be
the hallmark of the highest of the positive religions;
and the highest and concluding of them is Christianity.
Only by carrying through, as consistently as is at all possible,
an abstractly spiritualistic conception does
Christianity attain a distinctive view of life. All antecedent
stages, all preparations, are already given in the
pre-Christian religions; only that one thing remains over for it.

In the basic Christian representations the abstractly
spiritualistic is everywhere easy to demonstrate. It comes forth in the doctrine
of \emph{God} as the pure spirit who dwells in the light in which
there is no darkness of matter. --- It comes forth in the representation
of a \emph{creation out of nothing} by the almighty will:
for thereby the created is posited as that which in its essence is a
nothing, as that which has only from the arbitrariness of omnipotence, by
which it is upheld, a borrowed being that can at any
moment be annihilated, as happens partially in the miracle, and for
the whole of nature in the destruction of the world, in which what
% --- p. 450 ---
\opage{450}was a nothing returns to the nothing from which it
came, and lets only the spiritual remain. --- It comes
forth in the whole representation of the activity of the divine
as a \emph{miracle-working} activity, for in the miracle omnipotence works
without anything real; it denies the being of the real and its
validity. --- It comes forth in the doctrine of the \emph{human
task of life}, which becomes this: to seek unity with God
through Christ by deadness to the world, by the denial
of the nature that ``lies in wickedness'', which, when
it asserts itself as nature in us, is determined as the sinful,
and which therefore becomes that to which the spirit is to relate
only negatively, so that it does not become a material that
is to be formed, but one that is to be cast away, the
Christian being bound to flee the relations of natural life, which for him
have no significance but can only become defiling,
because they cannot be placed in a positive relation to the abstractly
spiritual (thus marriage, political life, property). ---
It comes forth in the conception of \emph{life after death} corresponding
to this conception of the task of life. What stands before
consciousness as an ideal, but as an ideal it
despairs of being able to reach in this existence, it transfers
to a life beyond, where everything natural that
still hinders and weighs down the spiritual shall have vanished,
where the spirit shall be freed from every natural limitation,
and where even the body in which imagination clothes the
departed spirit in order to be able to intuit it shall be a ``spiritual
body''. --- It comes forth finally in the very doctrine of the \emph{God-man},
precisely at that central point of Christianity
where the divine is united with the human and,
through the human, seems also to enter into a positive
relation to the natural. For this semblance vanishes as soon as
one considers that the life of the God-man is through and through
a \emph{miracle} and by the miracle is loosed from the usual natural conditions
of human life: at birth and in death,
and at every point between them. The unity of God and
% --- p. 451 ---
\opage{451}man in the God-man becomes only the unity of God and
the \emph{natureless} human. ---

Thus we have, by different paths, arrived at the
same result: that the character of Christianity is natureless
spirituality, which stamps its conception of the divine
and of the human. The consideration of the
historical presuppositions, of the conceptual development of the positive religions,
and of the basic doctrines of Christianity itself
has led us to one and the same conception and, by
this agreement, given it certainty. And the feeble attempts
that are made from the side of the theory of mediation to dispute this
conception are of such a kind that they will not easily be able to make anyone
doubtful. We will in passing --- for that will be
sufficient --- cast a glance at one such attempt, which
is found in the collection of edifying reflections on ethical
subjects that Bishop Martensen has recently published.

It begins with the postulate that Holy Scripture, on
every one of its pages, is supposed to bear witness that God is not the natureless
spirit. ``Scripture knows only the living God. But the
living God we (!) cannot think otherwise than in relation
to an eternal nature subordinate to him.'' It then continues:
``When Scripture speaks of God as an \emph{indissoluble}
life (Hebr. 7), we cannot think of this either except
as the indissoluble unity of the oppositions of life, among
which we know only the opposition between spirit and nature,
the ethical and the physical.'' This exegesis, which by its
mixture of childishness and hair-splitting recalls that of
the rabbis, is as unfortunate as possible, as anyone can
convince himself by reading through two lines at
% TR: p. 451 no full stop after „(V. 16 f.)“ in the print; the English supplies one.
the passage cited (v. 16 f.). Christ is compared
with Melchisedek, as the eternal high priest in contrast
to the mortal high priest under the law of Moses. As the one
who has no ``beginning of days or end of life'', as
the one who ``remains high priest continually'', he is called high priest
``after the power of an indissoluble life'', and this expression is interpreted
in what immediately follows, where it is said: ``\emph{for}
% --- p. 452 ---
\opage{452}thou art a priest for \emph{ever} after the order of Melchisedek.'' Here
there is, then, quite simply an opposition between the life that
can cease in death and that over which death has no
power; but not a single turn of phrase points toward that
higher unity of ``the oppositions of life'' which we ``must
think''. One must be uncommonly destitute of reasons to
resort to such an art of interpretation. Yet there is still
one argument left that almost surpasses the previous ones. It
% TR: p. 452 the print's „Anthropomofismer“, „Anthropomofismes“, „Anthropomofistiske“ (twice) rendered "anthropomorphism(s)", "anthropomorphic".
is drawn from the fact that Scripture everywhere speaks of God in anthropomorphisms.
``But the truth of religious anthropomorphism
rests on the nature in God.'' ``Living piety will not let
itself be disputed out of the belief that God has eyes and ears, hands
and feet, an arm that is not shortened; and it does not give up
seeing his finger in the events of the world and in the
life of the individual; for even though it is conscious of the figurative, the symbolic,
and recognizes that everything must be removed that
belongs only to created limitation, it nevertheless continues to
hold fast that there must be a real counterpart to all this
in God; which in other words means that God must
have organs of revelation, instruments --- an organism.'' ---
So: the anthropomorphic is to be image, symbol,
hence \emph{not} express actuality; but we are nevertheless to demand
a \emph{real} counterpart, and while we remove everything that could
give these determinations content, we are to ascribe to God a
body, an organism. It is the old method characterized
above: one gives with one hand and takes
with the other. The anthropomorphic is to be non-real
and yet real, non-literal and yet literal, and that in order
to support a poor, tottering, half-speculative theory that fables
about ``physical all-possibilities'' in God, about a ``relative nothing''
that is to be the same as the ``eternal, physical possibilities'', about
% TR: p. 452 only „over“ of „overnaturlige“ spaced; emphasis on "super".
``physical but \emph{super}natural powers'', about a matter that
differs from nature by having ``stepped out (how?)
of connection with the spirit'', and the like. It is characteristic
that the herald of the theory must himself concede that ``no
concrete comprehension'' of it can take place. Indeed! the
% --- p. 453 ---
\opage{453}empty phrases dissolve of themselves into a nothing that is not
a ``relative'' but an altogether complete ``absolute'' one, and it is
by an ``indefensible misunderstanding'' that the authors of Holy Scripture
are charged with professing such a theory.
What would the author of the Gospel of John say to the
interpretation cited of ``the figurative language of Scripture?''

\begin{center}\rule{1.8cm}{0.4pt}\end{center}

If we hold fast to the conception of Christianity as
that whose hallmark is natureless spirituality, then
the historically given phenomenon that meets us
under the name of Christianity becomes easily explicable to us. We understand
not only Christianity in its struggle with the world, in its
classical figures, in the forms of life in which it has sharply
stamped its distinctive essence; but we also understand the
forms in which the Christian is mixed with its opposite.
We understand without the least difficulty how the nature-denying
and world-denying religion --- as it
(contrary to the hope of the first Christian age for Christ's imminent
return) was to go on existing in the alien world,
as it was to be proclaimed in this world, as it was to
defend itself against the attacks from it, and as it was to be unfolded
in consciousness --- had to exchange its originally merely
negative relation to nature, to the natural faculties and
powers (naturally determined thought included), and to
the concrete forms of the life of the world, for a forced connection
with what it denied, which was taken up as means for the
Christian. Thus we explain what has been called a
Christian science, a Christian art, a Christian morality,
a Christian state. But we understand at the same time that the
forced factual recognition of what was denied in principle
could never become a full and true recognition, and that
therefore the forced connection could not become a real
union; for what was taken up remained, as mere means, without
essential connection with the Christian. Thus too
% --- p. 454 ---
\opage{454}the relation shows itself everywhere where the Christian becomes conscious
of itself, as we have demonstrated above (no. 4, p. 291)
with regard to ``Christian'' philosophy.

But while the phenomenon, from this standpoint which
posits the opposition between the Christian and the Human,
becomes easily explicable, the mediation that
speculative theory attempts on the basis of the conception, characterized
above, of Christianity's relation to the
natural, runs up against insoluble difficulties. If the mediation
is taken really seriously, then Christianity, which
clearly and distinctly proclaims itself as a new creation, that
is, as an absolutely new beginning, and which for centuries
has been conceived as such by its confessors, becomes something purely superfluous,
since what it is supposed to be in unity with the Human
has already been given prior to Christianity. And all
the great classical figures throughout the whole Christian
era must then be conceived as one-sided, limited individualities
to whom the proper truth of Christianity has not
dawned. --- If restrictions are made and reservations taken,
so that the mediation becomes only half --- and as a rule speculative
theology ventures to go only so far when
it has to pronounce on principles --- then both
the Human and the Christian are distorted, and both become meaningless.
The Human is distorted and becomes meaningless in that,
despite all appreciative phrases, it receives only a sham actuality
and a sham significance. The Christian is distorted
in that the elements it takes into itself from what is alien,
although they are as it were instantly sublimated, nevertheless cast a tinge
over the Christian that robs it of its specific distinctiveness,
make it something without principle and therefore meaningless.

\begin{center}\rule{1.8cm}{0.4pt}\end{center}

But if the conception of the Christian as
natureless spirituality stands firm, then it also stands firm that there
must continue to be an opposition between it and the
% --- p. 455 ---
\opage{455}Human, despite the significance the Christian may have had for
the development and purification of the latter. The Christian era has
been the purgatory of the human spirit, the time of purification of the Human.
Natureless, abstract spirituality with its negative infinity
has as it were burned sensuousness and finitude
out of the human essence, and thereby made it possible for the spirit,
as it immersed itself in the richness of its essence, to give this richness
the form of an infinity full of content, and thus to be able to live
an eternal life in the reconciliation of the spirit within actuality.
It has been able to break down the barriers that separated the peoples from
one another, and to strengthen the consciousness of a humanity, a consciousness which
antiquity had already reached. And Christianity,
by powerfully reminding the human spirit that human life
first finds its truth in the relation to God, has given
it the impulse as it were to deepen all the human relations of moral life
in the relation to God, to make them all religious
relations. But the view of life of natureless spirituality
includes consequences against which the human view of life
must defend itself. As the natural foundation of human life
becomes a nothing, and more precisely such a nothing as, when it
asserts itself as a something, as an actuality, is seen
as evil, then the foundation of concrete moral
life, of all the forms of life which, like marriage, society,
the state, rest on a natural basis, is taken away: only
the ascetic flight from the world, only a renunciation of the whole actual content
of moral life remains. And when
nature vanishes from the essence of the merely spiritual God, then
God's activity in the nature posited by an incomprehensible act of omnipotence
becomes the miracle-working activity that leaves no
room for laws of nature; and with the laws of nature fall the laws of thought,
and the validity of thought vanishes. Everywhere
thought meets the miracle as the higher truth, and at the central point
of Christianity, in the doctrine of the God-man, of the unity of the eternal,
all-embracing divine with the particular human being existing under
the limitation of the conditions of finitude,
it meets the miracle in its most pregnant form, and it can
% --- p. 456 ---
\opage{456}only assert itself and the human essence, whose fundamental determination
it is, by separating itself from what must annihilate
it. ---

But as the separation thus comes forth between the
human consciousness and the religion which, as the highest and concluding form
of the historically given positive religions, is
the legitimate representative of positive religion, the
human consciousness does not turn away from the religious in its ideal form,
freed from the limitation of the positive; rather it
finds this ideal religious as it immerses itself in itself,
as it penetrates to the ground of the human essence, to
the point where this essence is bound indissolubly to its divine
origin. How this relation between the Human
and the ideally religious is to be determined more precisely, we shall
seek to show in a following section, which will conclude these developments.

\begin{flushright}\textbf{Constans.}\end{flushright}
\begin{center}\rule{2.6cm}{0.4pt}\end{center}

\clearpage
% --- p. 492 ---
\opage{492}\begin{center}\textbf{\large The Christian and the Human as Opposites in the Consciousness of Our Time.}\end{center}

\secnum{III.}

When, in the many forms under which the human view of life
--- in its more clearly or more dimly conscious opposition
to the Christian --- has given itself expression since the
beginning of modern times, we seek what is common and pervasive,
what must be posited as the essentially determining factor, we
shall find it in two main points. The one is the striving of the
human orientation of life to assert the \emph{undisturbed, harmonious
unity in the human essence}; the other is its striving to
assert \emph{the ethical in its concretion}, to give the moral
life a content of actuality rich enough to fill out \emph{the whole}
of human life. But these two main points stand, as will readily
appear, in the closest inner connection.

Whereas the Christian view of life --- by denying the validity
and essentiality of the natural, and by placing truth in that
which for thought was the unappropriable, that which conflicted
with all the laws of cognition --- had set a rupture in the human
essence, a rupture that was infinitely sharpened when the
doctrine of original sin posited human nature as disturbed and
corrupted, the human view of life asserts unity in the
\emph{double} sense of a unity between the natural side and
% --- p. 493 ---
\opage{493}the spiritual side in man, and of a unity, a harmony, between
knowledge and faith, between cognition and will, between the
theoretical and the practical. The nature which the medieval
spirit had denied, from which it shrank back in dread, while
the imagination was irresistibly lured by what it dreaded ---
\emph{that} was restored to its rights once more at the beginning
of the new age. As the human spirit sought into itself in order
to find itself in its freedom and independence, it could no
longer remain blind to the indissoluble bond that bound it to
nature, to the fullness of life it drew from nature. The anxious
hovering between fear and allurement was replaced by enthusiastic
devotion. It was as if the Godhead had once more returned to the
forsaken nature, which was now again seen in essential connection
with its author, which in its beauty again became an object of
the spirit's admiration, and in its essentiality and its grounding
in law again became a source of cognition for the spirit. This
orientation toward nature runs through the whole of modern times.
We find it in the Renaissance's enthusiastic deification of the
natural universe; we find it in empirical science's devotion to
nature; we find it in the eighteenth century's rapture over nature
and in its demand for naturalness in education and in life; we find
it in our century in the renewal and spread of materialism,
theoretical and practical, in the mighty development of the natural
sciences, in philosophy's endeavor to see the unconscious and
conscious mental life in man as a unity and to give the divine a
natural content; indeed, we find it even in philosophizing theology's
groping attempts to get an imaginative picture of nature into the
supernatural God. But nature, once more acknowledged in its
essentiality, is then, as natural determinateness in the human,
seen in unity with the spiritual. Organic life unfolds into a
spiritual life, into a life of consciousness; from our natural
side, through the organic conditions, our cognition receives its
content; from our life of drives the life of will receives its
material. The natural in us becomes no
% --- p. 494 ---
\opage{494}longer that which is only to be denied, that from which our will
is to die away in order to live a purely spiritual life; it
becomes again, as in Greek antiquity, that which is to be formed
and harmonized by the will that, through cognition of our essence
and its fundamental relations, has gained clarity about its task,
and that realizes the demand of true cognition. In one-sided forms
this acknowledgment of the natural in us emerges in those
psychological, ethical and social theories that let our
determinateness by drives become the supreme ruling power, that let
it govern our individual development and the ordering of social
relations. But what asserts itself here in a one-sided way receives
its true form where an original lawfulness in our \emph{spiritual}
essence is acknowledged, according to which the natural
determinations in us are ordered and formed, and under which they
bow in service while being acknowledged in their
legitimacy.

When the unity of the \emph{natural side} and the \emph{spiritual side}
in us is asserted, the unity between \emph{knowledge} and \emph{faith},
between cognition and will, is asserted as well. It is the denial
of nature that robs thought of its foothold; it is the miracle, in
which the actuality of the real is annihilated in favor of the
incomprehensible working of supernatural omnipotence, that sets
thought an insurmountable barrier and robs its laws of their truth.
Where faith, as in positive religion, is faith in the wonder-working
God, in the God before whose will the laws of nature vanish, at
whose command nature itself vanishes, just as it came into being at
that command: \emph{there} faith and knowledge must form irreconcilable
opposites; and since faith, as the organ of the higher truth,
condemns knowledge to submission, to self-surrender, a fundamental
form of the human spirit, an essential determination in it, is here
denied. A rupture then enters the essence, an inner discord; for the
essential that is to be denied cannot, precisely as essential, be
annihilated. The relation can be otherwise when the natural is
asserted in its truth, when nature is no longer held up over the
abyss of nothingness by an omnipotence hovering above it, but when it,
% --- p. 495 ---
\opage{495}as bearing the life of the divine within itself, rests in its
laws determined by its essence. Here the essence and working of the
divine does not become something from which thought is thrown back,
and which can be grasped only by a faith that denies thought: it
becomes a content for the knowledge that, following the demand
lying in the human essence, ventures out confidently from the finite
in its multiplicity and limitation to the infinite. And when
knowledge and faith no longer stand as opposites, when the human
spirit finds rest in the content of its cognition, there comes to be
a harmony between cognition and action as well. The faith that stood
against knowledge also gave the will a law opposed to the law of
knowledge. Not in the demand of the human essence but in the command
of a supernatural, incomprehensible will was the norm of action
sought, where the task of action was not the purely negative one: to
flee the natural. But when faith and knowledge are in harmony, then
the demand of knowledge, which expresses the law of the human
essence, the human essence's own demand, is also seen as the one
valid demand upon the will; then the same undivided human essence
that comes to clarity about itself in its cognition is seen as
expressing itself, in agreement with itself, in conscious action.
The theoretical and the practical are then seen in harmony: one
essence unfolds through both. The spirit does not waver between a
double principle, a double law; it is not gnawed by doubt about the
truth of its cognition and the validity of its action: in the harmony
it has its assurance; it has rest in unity with itself. ---

As the \emph{second} main motive in the human orientation of life
we determined \emph{the assertion of the ethical in its concretion},
% TR: p. 495 only „Virkeligheds“ of „Virkelighedsindhold“ is spaced; the emphasis falls on "actuality"
of the moral life's content of \emph{actuality}. Within the religion
whose character was abstract, nature-denying spirituality, the
ethical task of life could consistently become only a dying away
from everything natural, an ascetic flight from the world in order
to find rest in a purely spiritual life. The ideal had to become a
life beyond, where all natural barriers had vanished, where
everything burdensome and defiling in the
% --- p. 496 ---
\opage{496}nature-determined earthly existence had been removed. As long as
man lived under the conditions of earthly life, he could not but
feel constantly, in his bondage to the natural, a restraining bond
and a seed of discord within himself. The promptings of drive in the
natural, through which nature asserted itself as a power in us, were
felt as temptation and defilement, and the gaze turned longingly
toward the future in which the spirit would live freed from this
restraint. The asceticism of Christian antiquity and of the Middle
Ages, monastic life and the hermits' flight from the world, were not
one-sidednesses and exaggerations but consistent expressions of the
spirit of the religion. Where the natural has lost its validity,
where natureless spirit is the highest, \emph{there} no other life can
consistently be lived than one that is at every moment a dying away
from the world. For the concrete forms of moral life there can be no
room; their validity is conditioned by the validity of the natural
foundation in us. All moral relations and all moral forms of society
point back, directly or indirectly, to this natural principle; they
fall when it is denied. But the denial of it is for man a
self-contradiction, because the human essence comprises the
spiritual and the natural in inseparable unity. The force by which
nature in us is to be subdued can come only from the natural ground
itself in us, which is an abiding constituent of our essence: in the
very negation of the natural, the natural affirms itself. The life of
natureless spirituality therefore remains always only an attempt,
which can be repeated again and again but never carried through,
because that which is to be killed draws strength, in the very
struggle, from the hidden ground. Therefore there is anxiety and
unrest in this life, and it can be removed only when everything that
belongs to the human essence is acknowledged, when this essence in
its whole compass becomes the object of ethical action, when the task
becomes the free realization of the human essence in its
completeness, when \emph{the whole} of life in its fullness of actuality
is drawn in under the ethical.

% --- p. 497 ---
\opage{497}A striving toward this we also find, in many forms, less perfect
and more perfect, throughout the whole of modern times. It is such a
striving that emerges in the secular morality of civic life, which at
the close of the Middle Ages emancipates itself from the
ecclesiastical; it is this striving that emerges when religion is
transformed into morals and in this transformation is often
trivialized; it is this striving that runs through the philanthropic
endeavors, through socialism's attempts to order social relations,
through philosophy's endeavors to work ethics into actuality, to
determine the principle of moral action in such a way that it aims at
the whole richness of life.

But since the striving of the human orientation of life to assert
moral life in its concrete actuality points back to the assertion of
the validity of the natural, it thereby points back to what we
singled out as the \emph{first} motive: the assertion of the unity of
the essence. It is the human essence that in its undivided unity
harmoniously joins nature and spirit which can unfold itself in a
concrete moral life. But the \emph{unity} of the human essence is
inseparable from the \emph{undisturbedness} of the human essence; and
since the striving of the human orientation of life aims at the
double goal we have indicated, which lies on one and the same line,
it is borne by faith in \emph{the undisturbable good in human
nature}, in human nature's \emph{essential completeness} beneath all
the imperfection that is inseparable from its development and its
struggles. The firm conviction that the human essence cannot, through
its deed --- which, as the expression of the essence's activity, is an
affirmation of the essence --- transform, negate this essence, is a
\emph{pervasive} characteristic of all forms of the human view of
life.

\begin{center}\rule{1.8cm}{0.4pt}\end{center}

Since the human orientation of life separates itself in principle
from positive religion, it cannot surprise us that it has frequently
led to a separation from the religious altogether, to
% --- p. 498 ---
\opage{498}a ``morality without religion'', indeed even to a passionate
struggle against religion. This, however, is neither a pervasive
feature of the many forms of the human orientation of life, nor is it
its natural consequence. In many of those forms, e.g.\ in socialist
theories, we find a striving, however unclear, to give the view of
life a religious completion, which here as a rule does not escape the
fantastic. In other forms we find as it were a remnant of the
positively religious, rationalized fragments of it. In some, finally,
and these we regard as the properly consistent ones, we find from the
beginning a striving to assert the unity of the Human and the
religious in such a way that the latter is seen as most deeply
grounded in the human essence and as determined in its form and its
content by it. And it will also be easy to show that the cognition of
man's essence in its completeness must of necessity lead to the
religious relation, and that it can be completed only in this.

If what is peculiar to human life, that which by a decisive
difference distinguishes it from the lower forms of life, from plant
life and animal life, is the determination of self-consciousness, the
free knowledge of oneself and a knowledge of an infinite, then this
points back to the infinity of the human essence. Only the being that
has inner infinity can, in its relation to itself, be freed from the
bondage to the particular that characterizes the animal's
self-sensation, which is never loosened from its momentary
determinateness; only such a being can become conscious of itself in
its unity and universality. And only the being that has inner
infinity can have consciousness of an infinite, for this presupposes
the infinity of consciousness, and this in turn the infinity of the
being whose consciousness has this determination. But the spiritual
being that has inner infinity can reach the goal of its unfolding,
can find rest and satisfaction, only in a relation of infinity, in
such a relation as that in which its determination of infinity is
realized, as the being unfolds itself as a spiritual life-unity in
which every original determinateness, every realized aptitude, is
made a moment, drawn in under
% --- p. 499 ---
\opage{499}the conscious harmony of the being, a harmony realized through
the activity of spirit that becomes the primitive expression of the
self-conscious being: through cognition. And this relation of
infinity, which is thus most immediately determined as a relation of
the spirit \emph{to itself}, cannot be thought through without its
unfolding into a relation \emph{to the divine}, into a \emph{religious}
relation. When we seek the \emph{source} of the inner infinity in the
human, when we think through the concept of self-consciousness, we
arrive of necessity at the divine. The human self that becomes
conscious of itself in its inner infinity can become conscious of
itself only as a \emph{derived} self, as a self that has its ground in
something higher; and thus self-consciousness itself encloses a
relation to the source of the human essence, to its divine ground,
which is at the same time the ground of the whole of existence. Human
self-consciousness is consummated only in God-consciousness. Because
the human essence is grounded in the divine, it can, as a
self-conscious being, come to realization only as this relation, by
which the being has being, comes to be for its consciousness, and
comes to be for it in such a way that the relation to God goes in
unity with the self's relation to itself in self-consciousness.

Thus the inner connection of the religious relation with the
essential determination of the human becomes clear. But in order that
this relation shall acquire the character of \emph{true} infinity,
shall be freed from the limitation of finitude that constantly remains
attached to the positive religions, with their deity separated from
the world, personified, wonder-working, and with their separation of
the religious relation from the human spirit's other essential
relations of life, it must become a relation to a divine that is in
truth an infinite, in that, as governing and encompassing existence,
it is \emph{indwelling} in it, and, as grounding the human spirit, it
is \emph{of one essence} with it. And it must become a relation that
can in actuality permeate \emph{all} the essential relations of the
spirit, that can find its expression in the relations of \emph{moral}
life, and that, while raising them into a higher sphere, itself gains
fullness and actuality
% --- p. 500 ---
\opage{500}through them. Thus the religious relation comes to comprise
\emph{the whole} of human life, \emph{all} the spirit's expressions of
life; it becomes that in which the spirit at every moment finds rest
and reconciliation: reconciliation in its own inner being through the
harmony of the essence, and reconciliation with the ground of its
existence.

The religious relation so determined, the \emph{ideal} form of
religiousness, the religion that can become \emph{humanity's} religion,
thus lies in man's essential determination as a necessary,
inescapable demand; this relation, taken in its full, comprehensive
significance, becomes the highest, completing form of realization of
human life.

Only one determination still remains to be expressly stressed, lest
what has been said here be misunderstood. The source of religion has
been sought in the human essence; religion has been seen as having
sprung from this essence, and from this point of view religion is
seen as the work of the human spirit. But since the human essence is
itself grounded in the divine, the origin of the source is to be
sought deeper: in the being through which the human spirit has
infinity as an inner essential determination, and the need for the
realization of this infinity. From this point of view religion then
becomes \emph{the work of the Godhead}, a work of its inner revelation
to the human spirit, i.e.\ of its \emph{communication of essence} to
it.

\begin{center}\rule{1.8cm}{0.4pt}\end{center}

Through the conception of the \emph{religious relation} in its ideal,
universally human form, the conception of \emph{the divine} to which
the human orientation of life must consistently lead has been
intimated. It cannot appropriate the positive religions'
representation of the personal, wonder-working God. In the
determination of personality, as it is made to count in the positive
religions and by theology, it can see only a finitizing
determination,
% TR: p. 500 print „Anthropomosisme“ (printer's error for „Anthropomorfisme“); rendered "anthropomorphism"
an anthropomorphism that is transferred to God precisely in order to
keep him separated, in representation, from the human and
% --- p. 501 ---
\opage{501}from existence, and to free his action from the law of rational
necessity and make room for the miracle. Instead of this
determination of the divine it must seek another and higher one,
which, while it expresses God's essence as \emph{spirit}, frees it
% TR: p. 501 the spacing runs through „Endeligheds-“ into „in-“ of „Endelighedsindskrænkningen“; the emphasis falls on "in finitude"
from confinement \emph{in finitude}; which, while it designates God
as the source of self-consciousnesses, frees him from the limitation
and relative abstraction of the determination of self-consciousness.
And by this determination it will then contradict the representation
of the \emph{wonder-working} God. The God who through the miracle breaks
through the laws of nature and overturns those of thought, the
thinking spirit cannot acknowledge without denying its essence, which
it does not have from itself but from the Godhead. It must demand a
God whose working does not contradict itself, as the working of the
wonder-working God does, in that it cancels what it has itself
posited and thus shows itself to be finite. It must demand a God whose
\emph{essence} finds its expression in the laws of reason that govern
existence, and who therefore cannot break these laws.

This God, acting according to the rational necessity of his essence,
it must then, in order to preserve his infinity, necessarily conceive
as \emph{indwelling} in the existence that has its ground in God, and,
as indwelling in it, as uniting it in his unity and thus \emph{ideally}
distinguishing himself from existence as it is unfolded in the form
of \emph{real} being. As indwelling (immanent) in existence, as the one
who permeates existence with his life and governs it from within, God
must be determined as standing in a relation of \emph{unity of essence}
to real existence and to the human spirit. From the essence of the
effect one must infer back to that of the cause, since the cause can
only be the cause of that whose essential determinations it bears
within itself. As author of nature God must not merely be a
spiritual, an ideal being, but must have the determinations of
reality in himself, must be \emph{ideal-real} being. But as author of
the nature that finds its completion in life and in spirit, God must
be ideal-real being in the sense that ideality is the primitive
determination, the fundamental determination
% --- p. 502 ---
\opage{502}in the divine, the one that comprehends and explains the determination of reality.
Through God's essence, then, nature in us is at once
asserted in its essentiality and spirit in its ideal dominion.
The oppositions of existence point back to one principle as to
their source; in the world of nature and in that of spirit one
divine reason everywhere asserts itself, working according to eternal,
inviolable laws, according to laws that have their guarantee
in the very essence of the Godhead, whose expression they are. The
lawfulness of existence is, as having proceeded from the indwelling eternal
reason, an eternally valid lawfulness. --- As the goal of existence
and as that which determines finite spirits and becomes the goal
and norm of their action, the divine is finally
to be seen under \emph{ethical} determinations, as the good; but this
ethical determination of it, which definitively designates it as
spirit is bound up in the closest way with its logical and
physical determinations; it becomes only their fuller and
richer expression, not a determination that frees the divine
from the eternal laws of thought and, as theology would have it,
makes room for miracle-working arbitrariness.

\begin{center}\rule{1.8cm}{0.4pt}\end{center}

Through the conception of the human essence and its relation
to the divine, a conception to which the human orientation of life
must consistently lead, \emph{the life-task for man} is determined.
It can only be the free and conscious actualization
of the human essence according to its complete determinateness.
But this general formula for the life-task must receive a
fuller expression. We can give it that by first
seeing the human self as that which is an indivisible unity
of spirit and nature, and as that in whose various fundamental activities
the unity of the essence must express itself in accord
with itself. The task then takes the more definite form that
man is to develop his personality in the inner harmony of his
fundamental activities and in inseparable unity of the spiritual
and the natural. Man thus receives his \emph{essence} as a
% --- p. 503 ---
\opage{503}task, as the \emph{concrete} unity in which the real
natural foundation is not to be denied or cowed, but taken up
into spirit and thoroughly formed by it, so that the will receives
content and fullness from its natural ground, while in turn that which
proceeds from the natural foundation receives a norm and a
rule in the ethical demand, which proceeds from the essence of spirit
with the stamp of the original (the a priori). This
ethical demand, which bears the stamp of the unconditioned, and which originally
comes forth in abstract form through the consciousness
of an obligation as such, of an unconditionally valid
``Thou shalt!'', is then, in order to be able to form thoroughly the natural
foundation of the personality, to be unfolded, through the cognition of man's
essence, into a richness in which it corresponds to the richness in actuality
of life and of action. The original consciousness
of the good, of duty, is in the actuality of action
to become a consciousness of this determinate good,
of this determinate duty, and during this unfolding it is to
retain its certainty, so that the whole of life through and through receives
the stamp of the clear and conscious ethical.

We are pointed toward the same goal also when we see
the general formula given above for the human life-task
from another point of view. That man is to make
his \emph{essence} a task for himself according to its many-sided and complete
determinateness means not only that he is to develop it
in harmonious unity of spirit and nature and of the fundamental activities of spirit;
but it means also that he is to develop
it, not as isolated and atomistic, but as \emph{a distinctive
member in a whole of existence}, as incorporated into
this existence and as subordinated to its principle.
Only thus does the self develop according to its completeness;
torn loose from the vital connection with the whole, it is
precisely the incomplete, just as the single organ in an
individual organism loses its integrity when it is separated
from the whole that sustains it. But as the self becomes
conscious of itself as incorporated into something more comprehensive and as subordinated
to something higher, its \emph{assertion of itself} in
% --- p. 504 ---
\opage{504}its distinctiveness becomes one with its \emph{devotion to the
higher} that it is to serve; in other words: its self-assertion
and self-unfolding receive the stamp of the \emph{ethical}, its willing
of itself becomes a willing of a higher demand, of the law of existence
and of the will of the Godhead. Thus self-development
loses every stamp of egoism; the unfolding of
individual distinctiveness is taken into the service of the whole and of
the principle; the life the individual lives for itself, it lives
for the whole of humanity and for the divine.
And as life at every point receives the stamp of the \emph{ethical},
it receives at the same time, since the self everywhere in existence, in all the
relative relations, meets the divine, the stamp of the \emph{religious};
the whole of life becomes at once ethical and religious;
the religious relation does not, as limited to particular
times or to particular domains of life, separate itself off from the other
essential relations of life, but comprehends, pervades and
elevates all the others.

In this view of life, which rests on the conviction
of the infinite worth of the human, in virtue of the indwelling of the divine
in the universal human essence, undisturbed in its ground,
and which leads to the recognition of the
human in all its essential forms of appearance and
to a universal love of humanity, of the human
\emph{as} human, the \emph{Human} is thus in inward
unity with the religious, which is seen as proceeding from
the ground of the human essence and as gaining content through
the ethical. It energetically points man toward an activity
in the world of actuality, but toward an activity that
in the work of actuality preserves the connection with the
higher. By reminding at every moment of the solidarity of the
human race, it reminds the single individual of the necessity,
in order to solve his individual task, of working toward the solution
of the social tasks. While the view of life
that denies nature only gathers individuals in the work
for their blessedness, but cannot unite them in a work
for the earthly conditions of life for a life of spirit,
% --- p. 505 ---
\opage{505}the human view of life gathers them for this work, and in this very
work it lets them work for their spiritual goal. Only
through it can the foundation be laid for \emph{an existence worthy
of human beings for all}, because only it recognizes in the right relation
the significance of the spiritual and of the natural,
the right of the individual and the demand of society, self-regard and
solidarity. ---

This religious-human view of life does not, however, form for itself
an idyllic picture of human life, a utopia, freed from
the imperfections of actuality. It recognizes, over against
the infinite demand of the moral, the distance of human life
from its ideal, its imperfection, its guilt, its
sinfulness. But in recognizing the painful distance
between ideal and actuality, it does not give up faith in
man, in his capacity to actualize the good, in
giving \emph{this} existence a content of eternity. When
positive religion teaches the disturbance of the human essence
by sin, a disturbance that only a divine miracle
can remove, that religiousness of humanity maintains that
% TR: p. 505 the Danish spaces only „-væsenet“ in „Menneske\emph{væsenet}“; the English carries the emphasis on "essence" alone.
the human \emph{essence} cannot have been transformed by an act
that has proceeded from the essence itself, any more than the essence of the organism
is transformed because its integrity is suspended
during illness. And just as little as it can concede
that the human essence is disturbed by sin, just as little can
it concede that nature should, through man, have been drawn
into the disturbance. It cannot do so, because it sees
the whole of existence as sustained by the indwelling principle.

When it recognizes evil, not merely as an imperfection,
a lack, but as an opposition to the good, posited by man's will,
it sees in the will that
is active toward evil a will that must be broken in the struggle
against the good, because it struggles against what is the very
\emph{essence} in the human, and must call forth a reaction
from it. And it sees at the same time in this will, which has the
semblance of false infinity, but which receives this semblance only from
that which it struggles against and strives to displace, something active,
% --- p. 506 ---
\opage{506}the direction of whose working can be so reversed that the
infinitely strained energy actualized in the service of evil can
become active for the good, acquire positive worth for it. Just as
in the domain of the physical the forces can be converted into one another,
yet lose nothing in this conversion, so here, in
the transformation of force, the sum of force is to remain unchanged.
But for the good to be able to be actualized in the individual, this
view of life too must recognize that individuals need
assistance. This assistance, however, it does not seek outside the
existence in which the individual has its actuality; it seeks
it in the other, supplementing individuals with whom the individual
joins together in love, in mutual communication of essence,
and, in the last instance, in the principle that is indwelling in the individual itself,
and that joins it together with the other individuals
in the unity of the essence. As individuals enter into a relation of love
to one another, it sees them as acting --- at
once purifying and supplementing --- upon one another. The
single individual becomes for the other a representative of
man's universal, reason-determined essence, and as such
it becomes for the individual, as it were, its personified, objectified
conscience. As the one individual in that relation
converts the other into a representative of the ideally human,
it gains in it a determining norm, an ideal standard
for its action, a corrective for it. And as the intuition
of the universally human fuses with
the intuition of the distinctiveness of personalities, of the distinctive
forms in which the ethical is actualized in the different
individuals, then the ethical distinctiveness, idealized by love, by the
unconscious faith in man, includes
not merely a corrective for what is deficient in the individual,
but also an assurance of the ethical power in
the human, and the individuals embraced by the relation of love
% TR: p. 506 the Danish breaks „Saa|ledes“ at the line end without a hyphen (a printer's defect, joined in the transcription); no effect on the English.
show themselves also as supplementing one another. Thus
the relation of love between human beings becomes an ethical
source of life for the individuals. But through this relation of love,
in which individuals are joined together in a spiritual
% --- p. 507 ---
\opage{507}unity, there is a pointing toward the principle of this unity, toward the origin of individual
spirits: the divine, and in it
the assistance for the individual is in the last instance to be sought. The
divine is, as the source of the individual spirit, that
which works through it; it is that from which
the human spirit has the determination of infinity that
comes forth in its relation to itself, in its grasping of
the good, and even in its egoistic willing of itself in opposition
to the good. From this its ground the human spirit can
never be separated, not even when it wills to be, for through
the evil will itself the divine will works with
the force of the essence, breaking and consuming the force of evil.
In the breakthrough in which the good will breaks forth
through the dissolution of evil, the divine is active;
for the essence asserts itself in the breakthrough only by a
working of \emph{that} which is the ground of the essence. And in every struggle,
in every wavering, the will can seek back to this ground,
which reveals itself in the innermost of the self through the ethical
demand, and it can find in it the completion of
the lacking force through the full self-surrender by
which the force of the eternal is won, as the individual freely loses
itself to the eternal. Here there no longer remains a doubt
about reconciliation as \emph{there}, where the individual, trembling, knows itself
dependent on the election or reprobation of an incomprehensible arbitrariness,
% TR: p. 507 the print has no full stop after „(Calvin)“ before the next sentence; the English supplies it.
on the ``dreadful decree'' of omnipotence (Calvin).
Here there is certainty about reconciliation, because the individual finds the
divine in itself as indwelling and becomes conscious of itself in its
unity of essence with it. But as the last condition
for the actualization of our life's goal thus becomes the divine
working in us, we are pointed back again to the religious
as the consummation of our ethical life. In this unity of
the ethical and the religious the former has its truth, the latter its
actuality.

\begin{center}\rule{1.8cm}{0.4pt}\end{center}

Since it has thus been shown that the human view of life
can be completed only in the religious, its normal
% --- p. 508 ---
\opage{508}relation to positive religion is thereby definitively determined. There must
be an \emph{opposition}, a clear and conscious opposition;
but what is essential in the positive religions can, despite \emph{it},
be preserved, and the positive, historical religions can be recognized
in their significance for the attainment of the \emph{ideal}
stage of religiousness that has been designated as \emph{the religion
of humanity}. It is only one-sided forms of the
human orientation of life that have fought in hostility against the
positively religious, and that have conceived it as a mere
distortion of the truth. As the human orientation of life
gains a deeper cognition of man's essence and of
his fundamental relations, it cannot avoid recognizing that
it is precisely the same striving for the actualization of the
human spirit's essence in a true relation of infinity, from
which the ideal form of religiousness proceeds, that also
has been what moved most deeply and what unfolded in the positive
religions, even though that striving, in being realized under conditions of finitude,
still shows itself in the positive religions as
burdened with a limitation of finitude that comes forth
in the conception of the divine, in the separation of the religious relation
from the other essential human relations, in the
forms of the human activities of spirit that come to
unfolding on the stage of positive religion, and in the relation
between them. Everywhere in the positive religions, therefore, there
can be shown, as it were, a refraction in consciousness: infinity
will momentarily pervade finitude, but
finitude will nonetheless remain standing as unsuspended and cling as something
burdensome and hindering to spirit's striving for infinity.

Only in the form of religiousness that unfolds from the
Human fully conscious of itself is the original striving of the human spirit
freed from the burdensome limitation of finitude.
The religious relation acquires the character of true infinity,
because it joins the human spirit with a divine that
is indwelling in it, and because it comprehends the whole of human life, all
the essential human relations, so that the religious can
become the principle of the totality of actual human life.
% --- p. 509 ---
\opage{509}In this religious relation, and only in it, can religious
reconciliation become complete; for only in it does spirit not have
over against it a divine being that, in the mystery of its incomprehensible
will, becomes what is alien to spirit, but a
divine in which spirit finds itself again and feels itself
confirmed; --- and only in it does spirit have full reconciliation
in its own interior, since all the essential activities of spirit
are in harmony, and since the life of religion is not lived by a
flight from the actuality of human life and by a denial
of essential forms of its expression. As the religiousness that
pervades the whole of human life, this religiousness of humanity
can be designated as the religion of \emph{actuality} and
as that which brings a reconciliation \emph{within actuality},
a reconciliation for the \emph{undivided} human essence. It
is the expression of the eternal spirit of religiousness as the life-principle
of the whole of human life. As religion it is in living
inner unity with science and morality. It has its
root in the deepest ground of the human essence, and is therefore, since
free and conscious spirit has found itself and found
this form of religiousness as the expression of its essence, something
unshakable and imperishable, independent of the changes of the times.

But precisely the cognition of this includes the cognition
of the truth that the positive religions contain. What is
most deeply grounded in the human essence must, throughout the whole
historical unfolding of human life, have made itself visible, even if
only momentarily and incompletely. Spirit's striving for infinity
must therefore also reveal itself through the positive
religions, and it will do so all the more, the more the
\emph{mystical} element in religion comes forward and reflection
bound by representation recedes. In the most immediate forms of religious
life, in the living, childlike, believing
conception, the religious consciousness in its obscure
longing comes nearest to the relation of infinity; where reflection gains
the upper hand, it moves away from it, and most where theology,
finite reflection on religion, becomes the representative
of the religious consciousness. Then the
% --- p. 510 ---
\opage{510}separation of finitude between the personal God and man and
nature, the opposition between God's essence and the human,
the opposition between God's miracle-working activity
and the laws of nature and of thought, is carried through as a matter of principle,
and with the relation of finitude, which through the separation and
the opposition fixes itself in consciousness, the discord is fixed,
% TR: p. 510 the print has „religtøse“, a printer's error (t for i) for „religiøse“; rendered "religious".
and the goal toward which the religious consciousness strives is
not reached. If theology then becomes half-speculative --- and
further than that it cannot reach without ceasing to be theology ---
then all that is attained by the mediation attempted between
faith and knowledge, between the antirational and the rational,
between the positively religious and the Human, is that consciousness
may become befogged and obscured, but there comes no
harmony, no actual unity. Then the thinking
spirit must seek unity with itself by its own way, and it
finds it by immersing itself in its essence and, as it
gains the complete concept of the Human, by seeing it include
the religious, which does not conflict with thought but confirms
thought, which does not flee from life but fills out, ennobles
and elevates life. Here, then, the goal is found, and from the goal
one can look back with recognition at the stages that lie
behind it; but the recognition of their relative truth
is inseparably united with the clear and free consciousness of
what separates them from the goal. And therefore there is
no striving, either, to preserve the semblance of the unity that in
actuality is not present: from the standpoint of free, human
religiousness, every neo-rationalist attempt,
after having stripped Christianity of everything that makes it
a positive religion, to claim for \emph{that} reduced, abstractly rational
sublimate of Christianity which one professes
the name of the Christian, must be unconditionally rejected as unprincipled and
confusing.

\begin{flushright}\textbf{Constans.}\end{flushright}
\begin{center}\rule{2.5cm}{0.4pt}\\[-10pt]\rule{2.5cm}{0.4pt}\end{center}

\end{document}
